% ============ CHAPTER 9 — RESULTS, DISCUSSION AND CONCLUSION ============
\chapter{Results, Discussion and Conclusion}\label{cha:results}

\section{Automated verification results}

The formal gate (latest full run, 2026-10-23 closure audit, as recorded in
\texttt{docs/porting/acceptance.md}):

\begin{longtable}{>{\raggedright}p{4.2cm} >{\raggedright\arraybackslash}p{7.2cm} >{\raggedright\arraybackslash}p{3.2cm}}
\caption{Final automated gate}\label{tab:finalgate}\\
\toprule
\textbf{Surface} & \textbf{Command} & \textbf{Result} \\
\midrule
Backend & isolated broker on 1884 +
\texttt{npm --prefix backend test} & \textbf{PASS}
--- 141/141 (camera ingest, near-animal, wrong-line, ACL and e2e
subsets) \\
Godot (42 suites) & \texttt{bash godot/tests/run\_tests.sh} & \textbf{PASS}
--- 42/42 including wrong-line visual/HUD, camera confirmation and
near-animal warning suites \\
ESP32 native logic & \texttt{pio test -e native} (Unity) & \textbf{PASS}
--- 16/16 \\
ESP32 firmware build & \texttt{pio run -e esp32dev} & \textbf{PASS}
--- RAM 15.5\%, Flash 87.1\% \\
ML service & \texttt{pytest ml/tests} & \textbf{PASS} --- 8/8 \\
Web build & \texttt{npm --prefix web run build} (\texttt{tsc --noEmit} +
Vite) & \textbf{PASS} \\
\bottomrule
\end{longtable}

\noindent Progression of the gate during the project (each step merged
green): backend 84/84 (Stage 2A) $\rightarrow$ 106/106 (Stage 2E) $\rightarrow$
109/109 (Stadium P1) $\rightarrow$ 118/118 (Stadium P2) $\rightarrow$
137/137 (Stadium P3) $\rightarrow$ \textbf{141/141} (closure audit); Godot
16/16 $\rightarrow$ 22/22 $\rightarrow$ 33/33 $\rightarrow$ 34/34 $\rightarrow$
36/36 $\rightarrow$ 39/39 $\rightarrow$ \textbf{42/42}; ESP32 native
14/14 $\rightarrow$ 15/15 $\rightarrow$ \textbf{16/16} (near-warning cases
added in Stadium P3).

\section{Live end-to-end scenario results}

\begin{longtable}{>{\raggedright}p{4.6cm} >{\raggedright\arraybackslash}p{10.4cm}}
\caption{Live recorded runs on the working stack}\label{tab:livescenarios}\\
\toprule
\textbf{Scenario} & \textbf{Observed behaviour} \\
\midrule
\endfirsthead
\caption[]{(continued)}\\
\toprule
\textbf{Scenario} & \textbf{Observed behaviour} \\
\midrule
\endhead
Wrong-line hold at B & Request \texttt{WL-T1-1} held T3 at configured hold
\texttt{holdPhysicalM=320}; T3 decelerated from 48\,km/h and settled at
physical 320.1\,m; released exactly when T1 completed clearance back on UP. \\
Meet with seed 4242 & \texttt{POST /api/meet \{seed: 4242, trigger: true\}}
armed the scheduler; a station window held T3 (\texttt{commandedSpeedKmh:
0}, 48$\rightarrow$0\,km/h) and recorded \texttt{HELD}; identical seed
reproduces identical history in \texttt{meets.test.js}. \\
Near $\rightarrow$ confirmed chain & Trackside UP-B0 observation produced
\texttt{SUSPECTED} (zone1 YELLOW, T1 target 30\,km/h, ESP32
\texttt{WILDLIFE\_NEAR}); two JPEG frames then confirmed
\texttt{CONFIRMED} elephant 0.92 (T1 target 0, ESP32
\texttt{WILDLIFE\_DETECTED}); the confirming frame was written to the
evidence cap; clearing left \texttt{RECOVERY\_REQUIRED} latched;
\texttt{RESET zone1} returned \texttt{NORMAL} without hazards, sent
\texttt{RESET} to ESP32, and T1 stayed at 0 until an
\texttt{/api/speed} 48 re-arm. \\
ESP32 cloud loop & Wokwi scenario over a Pinggy raw-TCP tunnel: serial
\texttt{MQTT connected}, then \texttt{STATE LOCAL WILDLIFE}; the live
twin observed the external detection (\texttt{external=true}) with
STOP/RED, followed by release while the backend stayed under
\texttt{RECOVERY\_REQUIRED}/RED. \\
Wokwi hardware-in-the-loop wildlife & Pressing and releasing the wildlife
input produced \texttt{STATE LOCAL WILDLIFE} on serial, Godot STOP/RED and
backend recovery behavior per the Stage-D checklist (PASS rows,
2026-10-21). \\
\bottomrule
\end{longtable}

\section{Visual verification}

All simulation figures in this report (Figures~\ref{fig:overview}
through~\ref{fig:cameraconfirm}) were regenerated for this report with the
repository's capture scripts under headless software rendering ---
the same method the acceptance gate uses for visual QA. In capture mode
the broker is intentionally absent, and the HUD honestly displays
\texttt{MQTT offline}; separately, the acceptance record documents the
live GPU-path runs (WSLg D3D12, $\approx$11.8\,ms/frame at 1280$\times$720
in a 150-frame MovieWriter run).

\section{Team contributions}

All four members of Team~14 contributed to every surface of the system ---
staged specification and planning, backend and broker security, the Godot
twin and its asset pipeline, the React dashboard, the ESP32 firmware and
its Wokwi verification path, the ML microservice, the automated test gates
and the acceptance documentation --- with the staged design documents,
plans and the acceptance record in the repository serving as the shared
record of each stage's review.

\section{Limitations}

Stated plainly, as the project does everywhere:

\begin{itemize}
  \item This is an \textbf{educational simulation}. It is not certified
        railway signalling equipment and must never be connected to real
        infrastructure, trains or operational systems. Stopping distance
        is a visualization estimate, not a braking-curve guarantee.
  \item \textbf{Physical hardware items were not run.} The Stage-D
        hardware checklist rows remain \emph{not run}; all ESP32 evidence
        is native-logic plus Wokwi cloud verification.
  \item \textbf{Real ML weights were not downloaded} in verification;
        most runs use \texttt{ML\_STUB=1} (deterministic elephant),
        and the real-weights path is implemented, tested at the
        input-validation level, but not exercised end-to-end in the
        recorded gates.
  \item \textbf{Capture-mode timing:} headless software-rendered frames
        are seconds apart, so the two-frame confirmation window is
        verified over live HTTP rather than consecutive captured frames.
\end{itemize}

\section{Conclusion}

The project's goal --- a \emph{runnable} digital twin of a wildlife-safe
railway line with a trusted control plane, role-aware dashboards, an
authenticated IoT bus, fail-safe edge hardware and staged
camera/ML-assisted wildlife protection --- is met within its declared
scope. The architecture principle (one authoritative backend, many
restricted clients; commands through REST; state through retained MQTT)
proved workable across five very different clients and survives every
failure mode the gate exercises: lost broker, stale clocks, replayed
envelopes, duplicate evidence, hazards behind the train, out-of-zone
commands and locally-debounced field inputs. Every subsystem in the
final gate is green, and the acceptance record ties each feature to
named suites and recorded live runs. The remaining gap --- real-hardware
qualification and real-weights ML verification --- is a matter of
exercise, not architecture.

\section{Future scope}

\begin{itemize}
  \item Run the Stage-D physical-hardware checklist on a real ESP32 with
        wired peripherals (the firmware and wiring are already
        validated in simulation and headers are hardware-ready).
  \item Exercise the ML service with cached model weights end to end
        (camera path recorded with the stub is architecture-complete).
  \item Extend wildlife sensing with additional modalities (multi-hut
        acoustic fusion, GSM-connected level crossings) and richer
        published telemetry retention for analytics.
  \item Scale the twin world to more lines/higher fleets (the
        backend's block model and the twin's track registry already
        generalise), and add multi-user concurrent operator sessions.
  \item Introduce proper certificate-based broker TLS in place of the
        current local-plain auth path for remote deployments.
\end{itemize}
